\clearpage\thispagestyle{empty}
\pdfbookmark[0]{Sources for the February 1805 Juno Report}{juno180502-sources}
\begin{center}\textbf{Sources for the February 1805 Juno Report}\end{center}
\noindent This AI translation follows Franz Xaver von Zach’s report as excerpted in \href{https://gdz.sub.uni-goettingen.de/id/PPN236007467}{Gauss’s \textit{Werke}, vol.~VI} (Göttingen, 1874), pp.~260--261. The corresponding passage of \textit{Ueber die Juno}, article XVII of the \textit{Monatliche Correspondenz}, vol.~XI, February 1805, runs from \href{https://www.digitale-sammlungen.de/en/view/bsb10538603?page=195}{p.~185} to \href{https://www.digitale-sammlungen.de/en/view/bsb10538603?page=198}{p.~188}.
\par\smallskip\noindent Carl Friedrich Gauss calculated Juno’s fourth set of elements using observations of his own and observations from Wilhelm Olbers. The star discussed in the report appears on \href{https://archive.org/download/histoireclestef00lalagoog/page/n162.jpg}{p.~135} of Jérôme de Lalande’s \textit{Histoire céleste française} (Paris, an IX [1801]).
\par\smallskip\noindent The collected editor’s note refers to the \hyperlink{native-111}{rules printed on p.~106} and to Gauss’s \hyperlink{native-124}{\textit{The Zodiac of Juno}}, published in the March issue on \href{https://www.digitale-sammlungen.de/en/view/bsb10538603?page=237}{pp.~225--228}.
\par\smallskip\noindent Translation and typesetting produced by AI. Journal scans from the Bayerische Staatsbibliothek.
\clearpage
